\chapter{الرموز}
\label{sec:notation}

يُرمز إلى أنواع العصبونات بالأحرف الأربعة الأولى من الاسم الإنجليزي لناقلها العصبي الرئيسي: \nt{GLUT}: الغلوتامات، \nt{GABA}: حمض غاما-أمينوبيوتيريك، \nt{GLYC}: الغلايسين، \nt{ACET}: الأسيتيل كولين، \nt{DOPA}: الدوبامين، \nt{SERO}: السيروتونين، \nt{HIST}: الهستامين، \nt{NORA}: النورأدرينالين، \nt{ADRE}: الأدرينالين، \nt{ASPA}: الأسبارتات (الجدول~\ref{tab:types}).

الحروف أقل من الكميات، والحرف نفسه يعني أحيانًا أشياء مختلفة في فصول مختلفة. في الجدول أدناه لكل معنى سطره؛ وفي العمود الأيسر الصيغة التي أُدخل فيها.

\begin{center}
\footnotesize
\setlength{\tabcolsep}{4pt}
\renewcommand{\arraystretch}{1.12}
\begin{longtable}{>{\raggedright}p{2.0cm}>{\raggedright}p{9.6cm}>{\raggedright\arraybackslash}p{2.4cm}}
\toprule
الرمز & المعنى & موضع إدخاله \\
\midrule
\endfirsthead
\toprule
الرمز & المعنى & موضع إدخاله \\
\midrule
\endhead
\bottomrule
\endfoot
$a$ & ميل خاصية النقل الخطية، $f(u)=au$ & \eqref{eq:feedback} \\
$a$ & مستوى \nt{ACET}: $0$ للقراءة، $1$ للكتابة & \eqref{eq:acetnet} \\
$a_i$, $a$ & تعب المجموعة: في مولّد الإيقاع؛ وفي أرجوحة النوم البطيء & \eqref{eq:halfcentre}, \eqref{eq:updown} \\
$a_S$, $a_P$ & حساسية إيقاع القلب للمسرّع وللمكبح & \eqref{eq:gasbrake} \\
$A(r)$, $A_0$ & استجابة المشبك للنبضة عند التردد $r$؛ وعند المشبك المستريح & \eqref{eq:depression} \\
$A_1$ & استجابة المستقبِل لحصة واحدة من الناقل & \eqref{eq:quantal} \\
$A$ & مساحة غشاء العصبون مع التغصنات & \eqref{eq:dendriteload} \\
$A_f$, $A_s$ & نسبة التصحيح التي تبقى حتى الحركة التالية، في العملية السريعة والبطيئة & \eqref{eq:tworate} \\
$b_i$ & التدفق الذاتي لناظمة الإيقاع & \eqref{eq:seroneuron} \\
$b$ & سرعة إتعاب العمل للمجموعة في مولّد الإيقاع أو في أرجوحة النوم & \eqref{eq:halfcentre}, \eqref{eq:updown} \\
$b$ & عدد البتات في عينة رقمنة واحدة & \eqref{eq:probedata} \\
$B_f$, $B_s$ & قوة تعلّم التصحيح من الخطأ، في العملية السريعة والبطيئة & \eqref{eq:tworate} \\
$c$ & تركيز الناقل في الحجم & \eqref{eq:seroneuron} \\
$c$ & معامل الارتباط بين ضجيج العصبونات المتجاورة & \eqref{eq:correlated} \\
$c$ & المعامل الذي يحرّك به فرقُ العدّين المضربَ & \eqref{eq:dishreadout} \\
$c_f$ & استجابة العصبون $C$: أكبر $s_{f,p}$ على كل المواضع & \eqref{eq:scstage} \\
$c_m$ & السعة النوعية للغشاء & \eqref{eq:dendriteload} \\
$C$ & كلفة الجهد: الفعل في زمن $\tau_a$ يكلّف $C/\tau_a$ & \eqref{eq:vigor} \\
$d'$ & قابلية التمييز بين دخلين: الفرق مقسومًا على الضجيج & \eqref{eq:gainsnr} \\
$d$, $d_j$ & التأخير على مسار الإشارة & \eqref{eq:glutneuron}, \eqref{eq:vstar} \\
$d$ & تأخير حلقة التغذية الراجعة مع العضلة & \eqref{eq:delaystable} \\
$D$ & معامل الانتشار & \eqref{eq:diffusionlength} \\
$D$ & المقام في ترددات شبكة \foreignlanguage{english}{\nt{GLUT}--\nt{GABA}} & \eqref{eq:isn} \\
$D$ & زمن انتظار المكافأة المؤجلة & \eqref{eq:patience} \\
$e$, $e_{ij}$ & أثر الأهلية في المشبك & \eqref{eq:threefactor} \\
$e$, $e_i$ & خطأ الخرج أو الحركة، الآتي على سلك الخطأ & \eqref{eq:lms}, \eqref{eq:adaptivefilter}, \eqref{eq:tworate} \\
$E_{\text{spk}}$ & طاقة نبضة واحدة مع عمل المشابك & \eqref{eq:energy} \\
$E$ & تردد مجموعة \nt{GLUT} & \eqref{eq:ei} \\
$E_E$ & جهد التوازن للمشبك المثير & \eqref{eq:shunt} \\
$f$, $F$ & خاصية النقل: التردد دالةً في الدخل & \eqref{eq:ratenet}, \eqref{eq:gain} \\
$f$, $f_0$ & تردد ضربات القلب؛ الإيقاع الذاتي لناظمة الإيقاع & \eqref{eq:gasbrake} \\
$f_s$ & تردد رقمنة قطب واحد & \eqref{eq:probedata} \\
$F(t)$, $F_1$ & قوة العضلة؛ قوة انتفاضة واحدة & \eqref{eq:forcesum} \\
$F_1$, $F_2$ & قوتا عضلتين تشدّان المفصل في اتجاهين متعاكسين & \eqref{eq:antagonists} \\
$g$ & حساسية المستقبل للتركيز & \eqref{eq:seroneuron} \\
$g_L$, $g_E$, $g_I$ & موصليات التسريب والإثارة والتثبيط & \eqref{eq:shunt} \\
$g$ & الكسب الذي يحدده \nt{NORA} & \eqref{eq:gain}, \eqref{eq:machine} \\
$g$ & التثبيط العام من \nt{GABA} في شبكة الذاكرة & \eqref{eq:acetnet} \\
$g$ & كسب بوابة الألم، المحدد من أعلى & \eqref{eq:gate} \\
$g$ & كسب مرحلة واحدة من الشلال الهرموني & \eqref{eq:cascade} \\
$g_j$ & المرشح $j$ بثابته الزمني الخاص عند دخل المرشح التكيفي & \eqref{eq:adaptivefilter} \\
$G$ & كسب حلقة التغذية الراجعة & \eqref{eq:feedback} \\
$G$ & سرعة التصحيح في حلقة ذات تأخير & \eqref{eq:delaystable} \\
$h$, $h_I$ & الدخل الخارجي لمجموعة \nt{GLUT} ولمجموعة \nt{GABA} & \eqref{eq:ei}, \eqref{eq:isn} \\
$h(t)$ & انتفاضة عضلية منفردة & \eqref{eq:twitch} \\
$H(t)$ & العتبة العليا لضغط النوم: ببلوغها تنام الآلة & \eqref{eq:sleeppressure} \\
$I$, $I_i$ & التدفق الخارجي إلى العصبون أو المجموعة & \eqref{eq:ratenet}, \eqref{eq:wta}, \eqref{eq:updown} \\
$I$ & تردد مجموعة \nt{GABA} & \eqref{eq:ei} \\
$I$ & سطوع المنبه أو شدته & \eqref{eq:photonnoise}, \eqref{eq:nakarushton} \\
$I$ & تيار الدخل الذي يشحن الغشاء & \eqref{eq:dendriteload} \\
$k$ & عدد النبضات في الرشقة & \eqref{eq:burst} \\
$k$ & كم مرة يجب أن تفوق الزيادةُ الضجيجَ & \eqref{eq:photonnoise} \\
$k$ & معامل التغذية الراجعة في الشلال الهرموني & \eqref{eq:cascade} \\
$k_s$ & القوة التي يضخ بها الألمُ التحسّسَ & \eqref{eq:sensitization} \\
$K$ & عدد سمات النتيجة & \eqref{eq:vectortd} \\
$K$ & صلابة المفصل & \eqref{eq:antagonists} \\
$K$ & كسب حلقة الشلال الهرموني، $K=g^3k$ & \eqref{eq:cascade}, \eqref{eq:cascadestable} \\
$K_\varphi$ & قوة مواءمة المولّد اليومي مع الإشارة الخارجية & \eqref{eq:pll} \\
$L$ & طول خط التأخير & \eqref{eq:itd} \\
$L$ & طول السلك من مستشعر الألم & \eqref{eq:paintime} \\
$L(t)$ & العتبة الدنيا لضغط النوم: ببلوغها تستيقظ الآلة & \eqref{eq:sleeppressure} \\
$M$ & مقياس الصلة السببية بين النذير والمكافأة & \eqref{eq:retro} \\
$M_k$ & توقُّع السمة $k$ & \eqref{eq:vectortd} \\
$M_{ik}$ & قوة التثبيط من عصبون \nt{GABA} رقم $k$ & \eqref{eq:machine} \\
$M$ & عزم الدوران في المفصل & \eqref{eq:antagonists} \\
$M$ & عدد خطوط الدخل إلى الخلّاطات & \eqref{eq:cover} \\
$n$, $\bar n$ & عدد النبضات في النافذة؛ ومتوسطه & \eqref{eq:rateerror} \\
$n$ & عدد مداخل العنصر العتبي & \eqref{eq:cover} \\
$n_\uparrow$, $n_\downarrow$ & عدد نبضات منطقتي ”أعلى“ و”أسفل“ في النافذة & \eqref{eq:dishreadout} \\
$N$ & عدد العصبونات & \eqref{eq:correlated}, \eqref{eq:energy} \\
$N$ & عدد أقطاب المسبار & \eqref{eq:probedata} \\
$N_s$ & عدد مواقع القذف بين عصبونين & \eqref{eq:quantal} \\
$p$ & احتمالية قذف الناقل لكل نبضة & \eqref{eq:burst} \\
$p$ & نسبة العصبونات النشطة في شبكة الذاكرة & \eqref{eq:acetnet} \\
$p$ & الاضطراب الذي يُتعلَّم تعويضه & \eqref{eq:tworate} \\
$P$ & القدرة المنفقة على الإشارات & \eqref{eq:energy}, \eqref{eq:dishpower} \\
$P$ & عدم اليقين في التقدير المخزَّن & \eqref{eq:kalman} \\
$P$ & فعالية ”المكبح“: أسلاك \nt{ACET} إلى القلب & \eqref{eq:gasbrake} \\
$P_{\max}$ & كم نمطًا عشوائيًا يستطيع العنصر العتبي أن يقسمها إلى صنفين & \eqref{eq:cover} \\
$P_{\text{miss}}$ & احتمالية الإخفاق & \eqref{eq:dishdensity} \\
$q$ & سرعة تغيّر العالم & \eqref{eq:kalman} \\
$q_k$ & تردد عصبون \nt{GABA} رقم $k$ & \eqref{eq:machine} \\
$q_0$ & الفعالية الخلفية للوسيط المثبِّط & \eqref{eq:gate} \\
$q$ & تردد الوسيط المثبِّط في بوابة الألم & \eqref{eq:gate} \\
$q$ & كم مرة تفوق الوحدة الحركية التالية سابقتها قوةً & \eqref{eq:sizeprinciple} \\
$r$, $r_i$ & تردد نبضات العصبون & \eqref{eq:ratenet} \\
$r$, $r_t$ & المكافأة (تدفق المكافأة) & \eqref{eq:rpe}, \eqref{eq:ctd} \\
$\mathbf r(t)$ & متجه استجابات الخزّان & \eqref{eq:reservoir} \\
$R$, $\bar R$ & المكافأة المتلقاة؛ وتوقُّعها أو متوسطها في وحدة الزمن & \eqref{eq:perturbation}, \eqref{eq:vigor} \\
$R$ & تباين الضجيج عند دخل المرشح & \eqref{eq:kalman} \\
$R$, $r_0$ & المكافأة المؤجلة والمكافأة الفورية & \eqref{eq:patience} \\
$R_{\max}$ & أقصى معدل للنقل، بت/ث & \eqref{eq:spikeentropy} \\
$r_{\max}$ & أعلى تردد للنبضات & \eqref{eq:gain}, \eqref{eq:nakarushton} \\
$R_{\text{raw}}$, $R_{\text{spk}}$ & تدفق بيانات المسبار: الخام والنبضات وحدها، بت/ث & \eqref{eq:probedata} \\
$s_j$ & علامة مخارج العصبون $j$ & \eqref{eq:dalesign} \\
$s_i$ & علامة مستقبل المستقبِل & \eqref{eq:seroneuron} \\
$s$, $\hat s$ & حالة العالم؛ وتقديرها & \eqref{eq:rpe}, \eqref{eq:kalman} \\
$s_{f,p}$ & استجابة العصبون $S$ للسمة $f$ في الموضع $p$ & \eqref{eq:scstage} \\
$S$ & فعالية ”المسرّع“: أسلاك \nt{NORA} إلى القلب & \eqref{eq:gasbrake} \\
$S$ & ضغط النوم & \eqref{eq:sleeppressure} \\
$t_1$ & لحظة النبضة الأولى & \eqref{eq:latency} \\
$t_k$ & لحظة النبضة رقم $k$ & \eqref{eq:forcesum} \\
$t_{f,i}$ & قالب السمة $f$ & \eqref{eq:scstage} \\
$T$ & نافذة عدّ النبضات؛ زمن تجميع الفوتونات & \eqref{eq:rateerror}, \eqref{eq:photonnoise} \\
$T_c$ & زمن صعود الانتفاضة العضلية & \eqref{eq:twitch} \\
$T_{\text{burst}}$ & الفاصل بين الرشقات الشبكية في المزرعة & \eqref{eq:burstinterval} \\
$u$ & الدخل الكلي للعصبون & \eqref{eq:gain} \\
$u$, $u(t)$ & أمر الدماغ: عند دخل المرشح التكيفي؛ وإلى الشلال الهرموني & \eqref{eq:adaptivefilter}, \eqref{eq:cascade} \\
$u(t)$ & دخل الخزّان & \eqref{eq:reservoir} \\
$U$ & مستوى الدخل الثابت & \eqref{eq:latency} \\
$U$ & نسبة مخزون الناقل التي تُقذف لكل نبضة & \eqref{eq:depression} \\
$v$ & سرعة الإشارة في السلك؛ سرعة حركة البقعة & \eqref{eq:itd}, \eqref{eq:vstar} \\
$V$ & الجهد على الغشاء & \eqref{eq:glutneuron}, \eqref{eq:shunt} \\
$V$ & توقُّع المكافأة المستقبلية & \eqref{eq:rpe}, \eqref{eq:ctd} \\
$w$, $w_{ij}$, $W_{ij}$ & وزن الرابط & \eqref{eq:glutneuron}, \eqref{eq:ratenet}, \eqref{eq:acetnet}, \eqref{eq:lms}, \eqref{eq:adaptivefilter} \\
$\mathbf w_k$ & أوزان مداخل العصبون $C$ رقم $k$ & \eqref{eq:traceinvariance} \\
$w_{q\mu}$, $w_{q\nu}$, $w_{z\mu}$ & أوزان الروابط في بوابة الألم & \eqref{eq:gate} \\
$W_{\text{out}}$ & أوزان القراءة الخطية للخزّان & \eqref{eq:reservoir} \\
$x$, $y$ & المداخل المنطقية والمخرج & \eqref{eq:dualrail}, \eqref{eq:veto} \\
$x$, $y$ & فعالية المرسل والمستقبِل & \eqref{eq:hebb} \\
$x$ & موضع الكاشف على خط التأخير & \eqref{eq:itd} \\
$x_i$ & الدخل من أعضاء الحس & \eqref{eq:acetnet} \\
$x_p$ & الدخل في الموضع $p$ & \eqref{eq:scstage} \\
$\mathbf x$ & مخارج العصبونات $S$، وهي دخل العصبون $C$ & \eqref{eq:traceinvariance} \\
$x_j$ & الدخل $j$ للمرشح التكيفي: الأمر بعد المرشح $g_j$ & \eqref{eq:lms}, \eqref{eq:adaptivefilter} \\
$x_f$, $x_s$ & تصحيحا العملية السريعة والبطيئة & \eqref{eq:tworate} \\
$x_1$, $x_2$, $x_3$ & المستويات في مراحل الشلال الهرموني الثلاث؛ $x_3$ الهرمون النهائي & \eqref{eq:cascade} \\
$x^*$ & نسبة مخزون الناقل المتعافي التي يبدأ عندها انهيار متسلسل جديد & \eqref{eq:burstinterval} \\
$y$ & دخل المرشح المشوَّش & \eqref{eq:kalman} \\
$y$, $\hat y$ & إشارة المستشعرات؛ وتنبؤ المرشح التكيفي بها & \eqref{eq:adaptivefilter} \\
$y_k$, $\bar y_k$ & فعالية العصبون $C$ رقم $k$؛ وأثرها & \eqref{eq:traceinvariance} \\
$y(t)$ & خرج قراءة الخزّان & \eqref{eq:reservoir} \\
$z$ & تردد خرج بوابة الألم & \eqref{eq:gate} \\
\midrule
$\alpha_i^\pm$ & وزن الأخطاء الموجبة والسالبة & \eqref{eq:asymmetric} \\
$\beta$ & قوة التثبيط المتبادل & \eqref{eq:wta} \\
$\beta$ & قوة تثبيط خرج بوابة الألم & \eqref{eq:gate} \\
$\gamma$ & معامل الخصم & \eqref{eq:rpe} \\
$\delta(\cdot)$ & دالة دلتا: قفزة لحظية & \eqref{eq:glutneuron} \\
$\delta$, $\delta_t$ & خطأ التنبؤ بالمكافأة & \eqref{eq:rpe}, \eqref{eq:ctd} \\
$\delta t$ & الفرق بين زمني وصول الصوت إلى الأذنين & \eqref{eq:itd} \\
$\Delta$, $\Delta_{\max}$ & الفاصل بين النبضات؛ نافذة التزامن & \eqref{eq:window} \\
$\Delta t$ & الدقة التي يميّز بها المستقبِل الزمن & \eqref{eq:spikeentropy} \\
$\Delta F$ & زيادة القوة من الوحدة الحركية التالية & \eqref{eq:sizeprinciple} \\
$\Delta I$ & الزيادة المميَّزة في السطوع & \eqref{eq:photonnoise} \\
$\Delta p$ & إزاحة المضرب في النافذة & \eqref{eq:dishreadout} \\
$\Delta u$ & الفرق بين دخلي المجموعتين & \eqref{eq:gainsnr} \\
$\Delta V$ & المقدار الذي يلزم رفع جهد الغشاء به & \eqref{eq:dendriteload} \\
$\Delta x$ & المسافة بين مستشعرين متجاورين & \eqref{eq:vstar} \\
$\varepsilon$ & الفرق النسبي بين الدخلين & \eqref{eq:wtagain} \\
$\eta$ & معدل التعلم & \eqref{eq:plasticity}, \eqref{eq:lms}, \eqref{eq:traceinvariance} \\
$\eta$ & الضجيج في الشبكة بعد المضخِّم & \eqref{eq:softmax} \\
$\mu$ & دخل اللمس إلى بوابة الألم & \eqref{eq:gate} \\
$\nu$ & دخل الألم & \eqref{eq:gate} \\
$\theta$ & عتبة الإطلاق & \eqref{eq:window}, \eqref{eq:scstage} \\
$\kappa$ & كمية الناقل لكل نبضة & \eqref{eq:seroneuron} \\
$\kappa_i$ & نسبة وزني الأخطاء الموجبة والسالبة & \eqref{eq:expectile} \\
$\kappa_m$ & صلة قوة العضلة بصلابتها & \eqref{eq:antagonists} \\
$\ell$ & المسافة التي ينتشر إليها الناقل & \eqref{eq:diffusionlength} \\
$\ell$ & ذراع العضلة التي تشدّ بها المفصل & \eqref{eq:antagonists} \\
$\xi$ & الارتجاف العشوائي لخرج العصبون & \eqref{eq:perturbation} \\
$\rho(\boldsymbol\psi)$ & نسبة الوقت الذي يُقضى في الضبط $\boldsymbol\psi$ & \eqref{eq:dishdensity} \\
$\sigma$ & التشتت، مقياس الضجيج & \eqref{eq:rateerror}, \eqref{eq:softmax} \\
$\sigma$ & السطوع الذي تبلغ عنده الاستجابة نصف أقصاها & \eqref{eq:nakarushton} \\
$\sigma_{\text{pre}}$, $\sigma_{\text{post}}$ & الضجيج قبل المضخِّم وبعده & \eqref{eq:gainsnr} \\
$\tau$ & الثابت الزمني للغشاء & \eqref{eq:glutneuron} \\
$\tau$ & الثابت الزمني لمرحلة الشلال الهرموني & \eqref{eq:cascade} \\
$\tau_{\text{eff}}$ & الثابت الزمني لمجموعة تثير نفسها & \eqref{eq:perfectint} \\
$\tau_c$ & عمر الناقل في الحجم & \eqref{eq:seroneuron} \\
$\tau_E$, $\tau_I$ & الثابتان الزمنيان لمجموعة \nt{GLUT} ومجموعة \nt{GABA} & \eqref{eq:ei} \\
$\tau_e$ & عمر أثر الأهلية & \eqref{eq:threefactor} \\
$\tau_{\text{tr}}$ & عمر أثر فعالية العصبون $C$ & \eqref{eq:traceinvariance} \\
$\tau_s$ & زمن عودة الحساسية بعد الإصابة & \eqref{eq:sensitization} \\
$\tau_r$, $\tau_d$ & زمن تراكم ضغط النوم في اليقظة وزمن تفريغه في النوم & \eqref{eq:sleeppressure} \\
$\tau_{\text{rec}}$ & زمن تعافي مخزون الناقل & \eqref{eq:depression} \\
$\tau_\gamma$ & أفق التخطيط & \eqref{eq:ctd} \\
$\tau_v$ & سرعة تحديث التوقُّع & \eqref{eq:ramp} \\
$\tau_a$ & زمن التعب في مولّد الإيقاع أو في أرجوحة النوم & \eqref{eq:halfcentre}, \eqref{eq:updown} \\
$\tau_a^*$ & أفضل زمن لأداء الفعل & \eqref{eq:vigor} \\
$\Phi$ & الطرف الأيمن لقاعدة التعلم & \eqref{eq:plasticity} \\
$\Phi$ & تدفق الفوتونات & \eqref{eq:photonnoise} \\
$\varphi_k$ & السمة $k$ للنتيجة المتلقاة & \eqref{eq:vectortd} \\
$\varphi$ & فرق الطور بين المولّد اليومي والإشارة الخارجية & \eqref{eq:pll} \\
$\boldsymbol\psi$ & ضبط الشبكة في نموذج \foreignlanguage{english}{DishBrain} & \eqref{eq:dishdensity} \\
$\omega$, $\omega_0$ & تردد الإشارة الخارجية (الضوء)؛ التردد الذاتي للمولّد اليومي & \eqref{eq:pll} \\
\end{longtable}
\end{center}
